\begin{theorem}[Unitary relabelling of transported vertical sectors]
    \label{thm:mpdo_transported_vertical_sector_unitary_relabeling}
    \lean{MPOTensor.transportedVerticalSector_exists_unitaryBlockEquiv}
    \leanok
    Under the hypotheses of Theorem~
    \ref{thm:mpdo_transported_vertical_sector_composites_identity}, write
    \begin{align}
      \mathcal A_1&=\prod_{\alpha}\MN{D^{(1)}_\alpha},
      &
      \mathcal A_2&=\prod_{\beta}\MN{D^{(2)}_\beta}.
      \notag
    \end{align}
    There are an equivalence $\sigma$ of the two sector index sets, equalities
    $D^{(1)}_\alpha=D^{(2)}_{\sigma(\alpha)}$, and unitaries
    $V_\alpha\in \MN{D^{(2)}_{\sigma(\alpha)}}$ such that
    \begin{align}
      \widetilde T(0,\ldots,0,X,0,\ldots,0)
      &=(0,\ldots,0,V_\alpha\,\iota_\alpha(X)\,V_\alpha^*,0,\ldots,0),
      \notag
    \end{align}
    where the nonzero entry on the right is in position $\sigma(\alpha)$, and
    \begin{align}
      \widetilde S(0,\ldots,0,Y,0,\ldots,0)
      &=(0,\ldots,0,\iota_\alpha^{-1}
        (V_\alpha^*YV_\alpha),0,\ldots,0),
      \notag
    \end{align}
    where the entries on the left and right are in positions $\sigma(\alpha)$
    and $\alpha$, respectively. These identities hold for every
    $X\in \MN{D^{(1)}_\alpha}$ and
    $Y\in \MN{D^{(2)}_{\sigma(\alpha)}}$.
    Here $\iota_\alpha$ is the reindexing induced by the displayed dimension
    equality. These are the mutually inverse unitary formulas of
    \cite[Appendix~C.4, line~1997]{Cirac2016MPDO_arXiv}. They make no assertion
    about the multiplicities, weights, or coefficient identity at
    lines~2001--2008.
\end{theorem}

\begin{proof}\leanok
    \uses{thm:direct_sum_inverse_cptp_forward_unitary,
      lem:cpsv_bnt_block_dimensions_positive,
      thm:mpdo_transported_vertical_sector_composites_trace_preserving,
      thm:mpdo_transported_vertical_sector_composites_schwarz,
      thm:mpdo_transported_vertical_sector_composites_identity}
    The two basis-of-normal-tensors decompositions have positive matrix
    dimensions. The transported maps are completely positive and preserve
    the total block trace, and the preceding theorem gives both inverse laws.
    Apply Theorem~\ref{thm:direct_sum_inverse_cptp_forward_unitary} to this
    pair. Its equivalence of summands, dimension equalities, and unitary
    matrices give the two displayed formulas without changing the chosen
    relabelling.
\end{proof}

\begin{theorem}[Coefficient comparison for transported vertical sectors]
    \label{thm:mpdo_transported_vertical_sector_coefficient_comparison}
    \lean{MPOTensor.transportedVerticalSector_exists_unitaryBlockEquiv_coefficient_eq}
    \leanok
    Under the hypotheses and with the notation of Theorem~
    \ref{thm:mpdo_transported_vertical_sector_unitary_relabeling}, let
    $\mu_\alpha$ and $\nu_\beta$ be the positive diagonal multiplicity
    matrices in the one-site and two-site vertical canonical forms, and set
    $m_\alpha=\tr(\mu_\alpha)$ and $n_\beta=\tr(\nu_\beta)$.
    For every horizontal bond matrix $X$ and every sector $\alpha$,
    \begin{align}
      A^{(2)}_{\sigma(\alpha)}(X)
      &=\frac{m_\alpha}{n_{\sigma(\alpha)}}
        V_\alpha
        \iota_\alpha(A^{(1)}_\alpha(X))V_\alpha^*.
      \notag
    \end{align}
    Consequently, for every tensor letter $a$,
    \begin{align}
      A^{(2)}_{\sigma(\alpha),a}
      &=\frac{m_\alpha}{n_{\sigma(\alpha)}}
        V_\alpha\,\iota_\alpha
          (A^{(1)}_{\alpha,a})V_\alpha^*.
      \notag
    \end{align}
    This is the coefficient identity in
    \cite[Appendix~C.4, lines~2001--2008]{Cirac2016MPDO_arXiv}. It compares
    the traces of the multiplicity matrices; it does not assert equality of
    their dimensions or of their individual diagonal entries.
\end{theorem}

\begin{proof}\leanok
    \uses{lem:mpdo_vertical_bond_contraction_encoded_matrix_unit,
      lem:mpdo_vertical_sector_retraction,
      thm:mpdo_transported_vertical_sector_unitary_relabeling,
      thm:mpdo_transported_vertical_sector_maps}
    Decompose the weighted contraction family into its individual matrix
    summands. At the summand paired with $\alpha$, the unitary formula gives
    \begin{align}
      m_\alpha V_\alpha
        \iota_\alpha(A^{(1)}_\alpha(X))V_\alpha^*
      &=n_{\sigma(\alpha)}A^{(2)}_{\sigma(\alpha)}(X).
      \notag
    \end{align}
    Positivity of $\nu_{\sigma(\alpha)}$ makes
    $n_{\sigma(\alpha)}$ nonzero, so division gives the first identity.
    Taking $X$ to be a matrix unit gives the tensor-letter identity.
\end{proof}

\begin{theorem}[Simultaneous representations of the blocked vertical operator]
    \label{thm:mpdo_blocked_vertical_operator_representations}
    \lean{MPOTensor.transportedVerticalSector_exists_blockedOperatorRepresentations}
    \lean{MPOTensor.verticalBNTMPO_verticalTensor_blockTwo}
    \lean{MPSTensor.sameMPV₂Pos_of_coisometry_reconstruction}
    \lean{MPOTensor.mpo_verticalBNTMPO_verticalAssembledTensor_eq_sum}
    \lean{MPOTensor.mpo_verticalBNTMPO_eq_sum_of_coisometry_reconstruction}
    \lean{MPOTensor.mpo_verticalBNTMPO_eq_pow_smul_of_unitary_reindex}
    \lean{MPOTensor.blockedVerticalOperatorRepresentations_of_unitaryBlockEquiv}
    \leanok
    \uses{thm:mpdo_transported_vertical_sector_coefficient_comparison,
      def:vertical_cf_mpdo, def:mpdo_two_site_blocking}
    Under the hypotheses and with the sector correspondence of Theorem~
    \ref{thm:mpdo_transported_vertical_sector_coefficient_comparison}, let
    $B$ be the vertical reading of the two-site blocking, and write
    $O_L(A)=\Tr_{\mathrm{bond}}(A_{i_1}\cdots A_{i_L})$
    for the length-$L$ closed-chain operator. If $L>0$, then the same
    equivalence $\sigma$, dimension identifications, and unitaries appearing
    in the tensor-letter identity satisfy
    \begin{align}
      O_L(B)
      &=\sum_\alpha
        \left(\frac{m_\alpha}{n_{\sigma(\alpha)}}\right)^L
        \tr\!\left(\nu_{\sigma(\alpha)}^L\right)
        O_L\!\left(A^{(1)}_\alpha\right),
      \notag\\
      O_L(B)
      &=\sum_{\alpha,\beta}
        \tr(\mu_\alpha^L)
        \tr(\mu_\beta^L)
        O_L\!\left(A^{(1)}_\alpha\right)
        O_L\!\left(A^{(1)}_\beta\right).
      \notag
    \end{align}
    These are the two representations in
    \cite[Appendix~C.4, lines~2011--2018]{Cirac2016MPDO_arXiv}.

    \textbf{Local fix (blocked coefficient exponent):} CPSV16 Appendix C.4,
    line~2013 prints $m_\gamma^L/n_\gamma$. The line-2008 tensor scaling and
    line~2040 give $(m_\gamma/n_\gamma)^L$. This is documented in
    \path{docs/paper-gaps/cpsv16_blocked_operator_trace_ratio_exponent.tex}.
\end{theorem}

\begin{proof}\leanok
    \uses{thm:mpdo_operator_product_mpo, lem:mpv_vertical_assembled,
      thm:mpdo_transported_vertical_sector_coefficient_comparison}
    The two-site vertical canonical form first gives
    $O_L(B)=\sum_\delta\tr(\nu_\delta^L)
    O_L(A^{(2)}_\delta)$.
    Reindex this sum by $\delta=\sigma(\alpha)$. The tensor-letter identity
    from the preceding theorem multiplies every letter of
    $A^{(1)}_\alpha$ by
    $m_\alpha/n_{\sigma(\alpha)}$. Along a chain of length $L$ this scalar
    therefore occurs to the power $L$; unitary conjugation and the bond-index
    identification leave the closing trace unchanged. This proves the first
    formula.

    For the second formula, the vertical reading of a two-site block is the
    product tensor of two copies of the one-site vertical reading. Closed
    chains turn this tensor product into multiplication of operators. Expand
    each copy by the one-site vertical canonical form,
    $O_L(\widetilde M)=\sum_\alpha
    \tr(\mu_\alpha^L)O_L(A^{(1)}_\alpha)$,
    and distribute the product of the two finite sums.
\end{proof}

\begin{theorem}[Canonical hypotheses for the retained vertical product]
    \label{thm:mpdo_retained_vertical_product_canonical_hypotheses}
    \lean{MPOTensor.verticalCoisometrySquare}
    \lean{MPOTensor.verticalCoisometrySquare_isCoisometry}
    \lean{MPOTensor.verticalTensor_blockTwo_squared_coisometry_reconstruction}
    \lean{MPSTensor.projectorClosure_and_noPeriodicVectors_of_coisometry_reconstruction}
    \lean{MPOTensor.retainedVerticalProduct_projectorClosure_and_noPeriodicVectors}
    \leanok
    \uses{thm:mpdo_blocked_vertical_operator_representations,
      def:mpdo_two_site_blocking}
    Suppose that the vertical reading of $M$ is reconstructed from a retained
    tensor $A$ by a coisometry $U$. The two-site blocking is then reconstructed
    from the product of two copies of $A$ by the Kronecker square of $U$.
    Moreover, if $M$ is in normalized BNT-refined horizontal form and generates
    positive operators, this retained product tensor has invariant-projector closure
    and has no nontrivial periodic vectors. No full-support assumption is
    imposed on the retained bond space.
\end{theorem}

\begin{proof}\leanok
    The Kronecker square of a coisometry is again a coisometry. Expanding a
    blocked vertical letter gives the product of the two one-site
    reconstructions and hence the asserted exact compression formula.
    Normalized BNT-refined horizontal form and positivity are preserved by
    two-site blocking. Invariant-projector closure descends to an exact reducing
    compression, while a periodic vector in the compression would embed
    isometrically as a periodic vector of the blocked vertical tensor.
\end{proof}

\begin{theorem}[Canonical hypotheses for every retained copy pair]
    \label{thm:mpdo_weighted_vertical_product_block_canonical_hypotheses}
    \lean{MPOTensor.productRetainedEquiv}
    \lean{MPOTensor.weightedVerticalProductBlock}
    \lean{MPOTensor.mulTensor_verticalAssembledTensor_reindex}
    \lean{MPOTensor.weightedVerticalProductBlock_projectorClosure_and_noPeriodicVectors}
    \leanok
    \uses{thm:mpdo_retained_vertical_product_canonical_hypotheses,
      lem:canonical_hypotheses_dependent_direct_summand}
    In the retained coordinates, the product tensor is the orthogonal direct
    sum over ordered pairs of BNT copies. Its summand indexed by
    $((\alpha,q),(\beta,r))$ is
    $\mu_{\alpha,q}\mu_{\beta,r}(M_\alpha M_\beta)$.
    Every such summand has invariant-projector closure and has no nontrivial
    periodic vectors. This includes zero-dimensional summands and requires
    no nonvanishing assumption on the displayed scalar.
\end{theorem}

\begin{proof}\leanok
    The canonical inclusion of a dependent direct-sum block is an isometry.
    After transporting this inclusion through the retained-coordinate
    equivalence, it intertwines the chosen summand with the full retained
    product tensor on both sides. Its range projection therefore commutes
    with every tensor letter. Invariant-projector closure descends by exact
    compression, and any periodic vector of the summand would embed as a
    periodic vector of the full retained product.
\end{proof}

\begin{theorem}[Normal-corner decomposition of a retained copy pair]
    \label{thm:mpdo_weighted_vertical_product_block_normal_decomposition}
    \lean{MPOTensor.exists_weightedVerticalProductBlock_normalDecomposition}
    \leanok
    \uses{thm:mpdo_weighted_vertical_product_block_canonical_hypotheses,
      thm:projector_closure_canonical_form}
    For every ordered pair of retained BNT copies, the weighted product tensor
    has an exact orthogonal decomposition into positive multiples of normal
    tensors. The corresponding corner isometries are retained, intertwine in
    both directions, and reconstruct every tensor letter. The active family
    is allowed to be empty; in particular, no normal block is introduced for
    a zero-dimensional or identically zero copy-pair tensor.
\end{theorem}

\begin{proof}\leanok
    Apply the canonical-form sufficient condition to the invariant-projector
    closure and periodic-exclusion properties of the preceding theorem. Its
    spectral construction discards precisely the corners on which every
    letter vanishes and retains the isometry of every nonzero corner.
\end{proof}

\begin{definition}[Simultaneous spectral data for retained copy pairs]
    \label{def:mpdo_retained_product_spectral_family}
    \lean{MPOTensor.retainedVerticalProductTensor}
    \lean{MPOTensor.retainedProductBlockInclusion}
    \lean{MPOTensor.RetainedProductSpectralFamily}
    \lean{MPOTensor.RetainedProductSpectralFamily.ActiveLabel}
    \lean{MPOTensor.RetainedProductSpectralFamily.activeLabelEquiv}
    \leanok
    \uses{thm:mpdo_weighted_vertical_product_block_normal_decomposition}
    A retained-product spectral family records, simultaneously for every
    ordered pair of retained BNT copies, all nonzero normal corners, their
    positive spectral coefficients, and their local isometries. Its active
    labels form the dependent union of the local corner families; this union
    is enumerated by a single finite type. The local isometries and their
    inclusions remain part of the data, and their range projections can be
    formed when needed.
\end{definition}

\begin{theorem}[Existence and inclusions of retained-product spectral data]
    \label{thm:mpdo_retained_product_spectral_family_exists}
    \lean{MPOTensor.verticalCopyCoordinateEquiv}
    \lean{MPOTensor.verticalCopyCoordinateEquiv_symm_apply}
    \lean{MPOTensor.verticalAssembledTensor_reindex_copyCoordinates}
    \lean{MPOTensor.verticalCopyBlockInclusion}
    \lean{MPOTensor.verticalCopyBlockInclusion_isometry}
    \lean{MPOTensor.verticalAssembledTensor_mul_verticalCopyBlockInclusion}
    \lean{MPOTensor.verticalCopyBlockInclusion_conjTranspose_mul_verticalAssembledTensor}
    \lean{MPOTensor.verticalCopyBlockInclusion_compression}
    \lean{MPOTensor.retainedProductBlockInclusion_isometry}
    \lean{MPOTensor.retainedVerticalProductTensor_mul_retainedProductBlockInclusion}
    \lean{MPOTensor.retainedProductBlockInclusion_conjTranspose_mul_retainedVerticalProductTensor}
    \lean{MPOTensor.exists_retainedProductSpectralFamily}
    \lean{MPOTensor.RetainedProductSpectralFamily.retainedInclusion}
    \lean{MPOTensor.RetainedProductSpectralFamily.ambientInclusion}
    \lean{MPOTensor.RetainedProductSpectralFamily.retainedInclusion_isometry}
    \lean{MPOTensor.RetainedProductSpectralFamily.retained_intertwine}
    \lean{MPOTensor.RetainedProductSpectralFamily.retained_intertwine_adjoint}
    \lean{MPOTensor.RetainedProductSpectralFamily.ambientInclusion_isometry}
    \lean{MPOTensor.RetainedProductSpectralFamily.ambient_intertwine}
    \lean{MPOTensor.RetainedProductSpectralFamily.ambient_intertwine_adjoint}
    \lean{MPOTensor.RetainedProductSpectralFamily.ambient_compression}
    \leanok
    \uses{def:mpdo_retained_product_spectral_family,
      thm:mpdo_weighted_vertical_product_block_normal_decomposition,
      thm:mpdo_retained_vertical_product_canonical_hypotheses,
      lem:matrix_sigma_block_inclusion}
    Suppose that the one-site vertical tensor has the stated exact
    coisometric reconstruction, that $M$ is in normalized BNT-refined
    horizontal form, and that $M$ generates positive operators. Then simultaneous retained-product
    spectral data exist. In copy coordinates, the assembled vertical tensor
    is the direct sum of the weighted simple blocks. The canonical inclusion
    of each copy is an isometry, intertwines in both directions, and its
    compression selects exactly that weighted block. Likewise, the canonical
    inclusion of a copy pair, followed by its local corner isometry, is an
    isometry into the retained product bond space and intertwines in both
    directions. Composing once more with the adjoint of the squared vertical
    coisometry gives an isometry into the bond space of the blocked vertical
    tensor. Under the exact coisometric reconstruction, these ambient
    inclusions intertwine each active weighted corner with the blocked
    vertical tensor in both directions. In particular, if $J_j$ is the
    ambient inclusion of an active corner, then
    $J_j^\dagger\mathcal V(M^{(2)})^iJ_j=\lambda_jA_j^i$.
    These explicit composite inclusions retain the range projections required
    for the sector-weight comparison.
\end{theorem}

\begin{proof}\leanok
    In copy coordinates the assembled tensor is block diagonal. The
    canonical dependent-sum inclusion is an isometry, and the two
    block-diagonal intertwining identities give its forward, adjoint, and
    compression formulas.
    Choose the exact normal-corner decomposition of every copy-pair tensor.
    The canonical copy-pair inclusions are isometries and select the weighted
    diagonal blocks of the retained product tensor. Composition with each
    local isometry gives the retained inclusion and both intertwining
    identities. Finally, write the blocked vertical tensor as
    $T=W^\dagger CW$ and the ambient inclusion as $J=W^\dagger R$. The
    coisometry identity $WW^\dagger=1$ transports both retained intertwining
    identities from $C$ to $T$; the compression identity then follows from
    $J^\dagger J=1$.
    For any blocked BNT label, choose its first multiplicity copy. Composing
    its canonical inclusion with the adjoint reconstruction coisometry gives
    an isometric reference corner of the blocked vertical tensor. Transport
    along an equality of bond dimensions preserves its isometry and exact
    compression formula.
\end{proof}

\begin{theorem}[Retained-product spectra from literal CPSV canonical form]
    \label{thm:mpdo_literal_cpsv_retained_product_spectral_family}
    \lean{MPSTensor.IsCPSVCanonicalForm.exists_retainedProductSpectralFamily}
    \leanok
    \uses{def:mpdo, def:cpsv_canonical_form,
      def:mpdo_retained_product_spectral_family,
      thm:mpdo_literal_cpsv_positive_blocking,
      thm:mpdo_vertical_projector_closure_literal_cpsv,
      thm:mpdo_vertical_no_periodic_vectors_literal_cpsv}
    Let $M$ generate matrix product density operators, and suppose that its
    doubled-index MPS tensor is in literal CPSV canonical form. Given an exact
    coisometric reconstruction of the one-site vertical tensor by a retained
    direct sum of weighted normal tensors, simultaneous retained-product
    spectral families exist for all ordered pairs of retained copies. Empty
    active families are preserved. This is the retained-product decomposition
    used in \cite[Appendix~C.4, lines~2020--2029]{Cirac2016MPDO_arXiv}.
\end{theorem}

\begin{proof}\leanok
    Literal CPSV canonical form passes to the two-site blocking. Proposition
    4.13 gives invariant-projector closure and excludes nontrivial periodic
    vectors for its vertical tensor. Apply the retained-product construction
    to these two properties and the squared coisometric reconstruction.
\end{proof}

\begin{theorem}[Flattened reconstruction of the retained product]
    \label{thm:mpdo_retained_product_flat_reconstruction}
    \lean{MPOTensor.retainedVerticalProductTensor_eq_sum_pairBlocks}
    \lean{MPOTensor.retainedProductBlockInclusion_orthogonal_of_ne}
    \lean{MPOTensor.RetainedProductSpectralFamily.retainedInclusion_orthogonal_of_ne}
    \lean{MPOTensor.RetainedProductSpectralFamily.retained_compression}
    \lean{MPOTensor.RetainedProductSpectralFamily.retained_reconstruction_active}
    \lean{MPOTensor.RetainedProductSpectralFamily.flatDim}
    \lean{MPOTensor.RetainedProductSpectralFamily.flatCoefficient}
    \lean{MPOTensor.RetainedProductSpectralFamily.flatBlock}
    \lean{MPOTensor.RetainedProductSpectralFamily.flatInclusion}
    \lean{MPOTensor.RetainedProductSpectralFamily.flatDim_pos}
    \lean{MPOTensor.RetainedProductSpectralFamily.flatCoefficient_pos}
    \lean{MPOTensor.RetainedProductSpectralFamily.flatCoefficient_ne}
    \lean{MPOTensor.RetainedProductSpectralFamily.flatBlock_normal}
    \lean{MPOTensor.RetainedProductSpectralFamily.flatInclusion_isometry}
    \lean{MPOTensor.RetainedProductSpectralFamily.flatInclusion_orthogonal_of_ne}
    \lean{MPOTensor.RetainedProductSpectralFamily.flat_intertwine}
    \lean{MPOTensor.RetainedProductSpectralFamily.flat_intertwine_adjoint}
    \lean{MPOTensor.RetainedProductSpectralFamily.flat_compression}
    \lean{MPOTensor.RetainedProductSpectralFamily.flat_reconstruction}
    \lean{MPOTensor.RetainedProductSpectralFamily.flat_sameMPV₂Pos}
    \leanok
    \uses{def:mpdo_retained_product_spectral_family,
      lem:matrix_sigma_block_inclusion,
      lem:same_mpv_exact_isometric_decomposition}
    Enumerate the dependent family of active corners by a single finite label
    set. The resulting coefficients are positive, the blocks are normal, and
    the inclusion ranges are pairwise orthogonal. For every tensor letter,
    \begin{align}
      C^i
      &=\sum_j W_j(\lambda_j A_j^i)W_j^\dagger,
      \notag\\
      W_j^\dagger W_l
      &=\delta_{jl},
      \notag\\
      C^iW_j
      &=W_j(\lambda_jA_j^i),
      \notag\\
      W_j^\dagger C^i
      &=(\lambda_jA_j^i)W_j^\dagger.
      \notag
    \end{align}
    Each compression $W_j^\dagger C^iW_j$ is
    $\lambda_jA_j^i$, and $C$ has the same positive-length closed chains as
    the finite direct sum $\bigoplus_j\lambda_jA_j$.
\end{theorem}

\begin{proof}\leanok
    First reconstruct the retained product as the sum of its canonical
    copy-pair blocks. Substitute the exact local spectral reconstruction in
    every pair, distribute the two inclusion maps across the finite sums, and
    identify the iterated sum with the dependent active-label sum. Finally
    transport this sum through the active-label enumeration. Isometry,
    orthogonality, intertwining, and compression are inherited by composition;
    the closed-chain equality follows from the exact isometric reconstruction.
    This is the active-corner decomposition in
    \cite[Appendix~C.4, lines~2020--2029]{Cirac2016MPDO_arXiv}.
\end{proof}

\begin{theorem}[Coverage of active product corners by the blocked BNT]
    \label{thm:mpdo_active_product_corner_blocked_bnt_coverage}
    \lean{MPOTensor.RetainedProductSpectralFamily.%
exists_blockedBNT_gaugePhase_of_flatBlock}
    \lean{MPOTensor.RetainedProductSpectralFamily.FlatBlockedBNTComparison}
    \lean{MPOTensor.RetainedProductSpectralFamily.%
exists_flatBlockedBNTComparison}
    \leanok
    \uses{thm:mpdo_retained_product_flat_reconstruction,
      thm:mpdo_retained_vertical_product_canonical_hypotheses,
      thm:cpsv_bnt_characterization}
    Suppose that the one-site vertical tensor has the stated exact
    coisometric reconstruction, and let $(A^{(2)}_\gamma)_\gamma$ be a basis
    of normal tensors for the blocked vertical tensor. For every active
    product corner $A_j$ there are a label $\gamma(j)$, an equality of bond
    dimensions, an invertible matrix $X_j$, and a scalar $\zeta_j$ of modulus
    one such that, after transporting the blocked tensor along the dimension
    equality,
    \begin{align}
      A_j^i
      &=\zeta_jX_jA^{(2),i}_{\gamma(j)}X_j^{-1}.
      \notag
    \end{align}
    The choices may be made simultaneously. No converse coverage or
    surjectivity of $j\mapsto\gamma(j)$ is asserted.

    \textbf{Scope restriction (active product BNT):} The conclusion gives
    only one-way coverage of active product corners. It does not assert that
    every blocked BNT label occurs. This is documented in
    \path{docs/paper-gaps/cpsv16_bnt_uniqueness_zero_coefficient.tex}.

    \textbf{Local fix (per-pair support):} Appendix~C.4, lines~2011--2029 of
    \cite{Cirac2016MPDO_arXiv} compares a sum over all retained copy pairs and
    does not isolate a fixed pair. Here fixed-pair support follows from the
    exact squared reconstruction and the canonical direct-sum corner; a zero
    pair contributes an empty active family. See
    \path{docs/paper-gaps/cpsv16_figure11_per_pair_support.tex}.
\end{theorem}

\begin{proof}\leanok
    The squared coisometric reconstruction gives equality of all
    positive-length closed chains between the blocked vertical tensor and the
    retained product. Transport the blocked basis of normal tensors along
    this equality. Proposition~2.7, applied to the flattened positive normal
    decomposition, then covers each active corner by a blocked BNT tensor up
    to an invertible gauge and a unit-modulus phase. Classical choice makes
    these comparisons simultaneous.
\end{proof}

\begin{theorem}[Distinguished blocked reference corners]
    \label{thm:mpdo_blocked_reference_corners}
    \lean{MPOTensor.blockedReferenceInclusion}
    \lean{MPOTensor.blockedReferenceInclusion_isometry}
    \lean{MPOTensor.blockedReference_intertwine}
    \lean{MPOTensor.blockedReference_intertwine_adjoint}
    \lean{MPOTensor.blockedReference_compression}
    \lean{MPOTensor.RetainedProductSpectralFamily.FlatBlockedBNTComparison.referenceInclusion}
    \lean{MPOTensor.RetainedProductSpectralFamily.FlatBlockedBNTComparison.%
referenceInclusion_isometry}
    \lean{MPOTensor.RetainedProductSpectralFamily.FlatBlockedBNTComparison.reference_compression}
    \leanok
    \uses{thm:mpdo_active_product_corner_blocked_bnt_coverage,
      lem:matrix_sigma_block_inclusion}
    Assume that every blocked BNT label has positive multiplicity and that the
    blocked vertical tensor has the exact coisometric reconstruction
    $T=U_2^\dagger C_2U_2$. Let $E_{\gamma,0}$ be the canonical inclusion of
    the first copy of label $\gamma$, and put
    $F_\gamma=U_2^\dagger E_{\gamma,0}$. Then $F_\gamma$ is an isometry,
    intertwines $T$ with the distinguished weighted BNT copy in both
    directions, and compresses $T$ to that copy exactly. For an active
    product corner covered by $\gamma$, transport $F_\gamma$ along the stated
    equality of bond dimensions. The transported map remains an isometry and
    obeys the corresponding transported compression identity.
\end{theorem}

\begin{proof}\leanok
    The copy inclusion is isometric and selects one diagonal block of the
    assembled tensor. Inserting $T=U_2^\dagger C_2U_2$ and
    $U_2U_2^\dagger=1$ gives both intertwining identities and compression.
    Transport along equality of bond dimensions preserves these equations.
\end{proof}

\begin{theorem}[Active product positivity from sector separation]
    \label{thm:mpdo_active_product_coefficient_phase_pos_of_sector_separation}
    \lean{MPOTensor.RetainedProductSpectralFamily.FlatBlockedBNTComparison.%
activeCoefficient_mul_phase_pos_of_sectorCompression_separation}
    \leanok
    \uses{def:mpdo, thm:mpdo_active_product_corner_blocked_bnt_coverage,
      thm:mpdo_blocked_reference_corners,
      thm:mpdo_retained_product_spectral_family_exists,
      thm:mpdo_sector_weight_pos}
    Assume that $M$ generates matrix product density operators and that the
    one-site and blocked vertical tensors admit exact coisometric
    reconstructions. Assume positive multiplicities and weights for the
    blocked BNT and normality of its representatives. Suppose moreover that
    every projection with a nonzero two-sided corner of the blocked vertical tensor
    has a nonzero finite-chain sector compression. Then every active product
    corner $j$ satisfies $0<\lambda_j\zeta_j$.
\end{theorem}

\begin{proof}\leanok
    Compare the active corner with the distinguished positive-weight copy of
    its blocked BNT representative. The active coefficient is nonzero because
    $\lambda_j>0$ and $|\zeta_j|=1$. The assumed separation property supplies
    a nonzero finite-chain compression, so the sector-weight comparison gives
    $\lambda_j\zeta_j>0$.
\end{proof}

\begin{theorem}[Positivity of active product coefficients]
    \label{thm:mpdo_active_product_coefficient_phase_pos}
    \lean{MPOTensor.RetainedProductSpectralFamily.FlatBlockedBNTComparison.%
activeCoefficient_mul_phase_pos}
    \leanok
    \uses{thm:mpdo_active_product_corner_blocked_bnt_coverage,
      thm:mpdo_blocked_reference_corners,
      thm:mpdo_retained_product_spectral_family_exists,
      thm:mpdo_grouped_sector_compression_nonzero,
      thm:mpdo_sector_weight_pos,
      thm:mpdo_active_product_coefficient_phase_pos_of_sector_separation}
    Suppose that $M$ is in normalized BNT-refined horizontal form and generates
    positive operators. Assume exact coisometric reconstructions of its one-site and
    blocked vertical tensors, positive multiplicities and weights for the
    blocked BNT, and normality of its BNT representatives. For every active
    product corner $j$, one has $0<\lambda_j\zeta_j$.
\end{theorem}

\begin{proof}\leanok
    Normalized BNT-refined horizontal form separates every nonzero
    range-projection corner at some finite chain length. Apply the preceding
    theorem with this separation property.
\end{proof}

\begin{theorem}[Literal positivity of active product coefficients]
    \label{thm:mpdo_literal_cpsv_active_product_coefficient_phase_pos}
    \lean{MPOTensor.RetainedProductSpectralFamily.FlatBlockedBNTComparison.%
activeCoefficient_mul_phase_pos_of_cpsvCanonicalForm}
    \leanok
    \uses{def:mpdo, def:cpsv_canonical_form,
      thm:mpdo_active_product_corner_blocked_bnt_coverage,
      thm:mpdo_blocked_reference_corners,
      thm:cpsv_literal_sector_compression_nonzero,
      thm:mpdo_sector_weight_pos,
      thm:mpdo_active_product_coefficient_phase_pos_of_sector_separation}
    Suppose that the doubled-index MPS tensor of $M$ is in literal CPSV
    canonical form and that $M$ generates matrix product density operators.
    Assume exact coisometric reconstructions of the one-site and blocked
    vertical tensors, positive multiplicities and weights for the blocked
    BNT, and normality of its representatives. For every active product
    corner $j$, one has $0<\lambda_j\zeta_j$. This is the positivity step of
    \cite[Appendix~C.4, lines~2025--2029]{Cirac2016MPDO_arXiv}.
\end{theorem}

\begin{proof}\leanok
    The literal sector-compression separation theorem supplies the finite
    chain on which the active corner is nonzero. Apply
    Theorem~\ref{thm:mpdo_active_product_coefficient_phase_pos_of_sector_separation}.
\end{proof}

\begin{theorem}[Unitary normalization of active product gauges]
    \label{thm:mpdo_active_product_gauge_unitary_normalization}
    \lean{MPOTensor.RetainedProductSpectralFamily.FlatBlockedBNTComparison.%
exists_unitaryNormalization}
    \leanok
    \uses{thm:mpdo_active_product_corner_blocked_bnt_coverage,
      thm:mpdo_retained_product_spectral_family_exists,
      lem:mpdo_horizontal_cf_block_two,
      thm:mpdo_block_two_as_positive_blocking,
      thm:mpdo_active_product_coefficient_phase_pos,
      thm:mpdo_blocked_reference_corners,
      thm:mpdo_pairwise_vertical_gram_conjugation,
      thm:normal_tensor_gram_rigidity}
    Under the hypotheses of the preceding positivity theorem, for every
    active product corner $j$ there is a real number $\omega_j>0$ such that
    \begin{align}
      X_j^\dagger X_j
      &=\omega_j\Id,
      \notag\\
      \omega_j^{-1/2}X_j
      &\in\mathrm{U}(D_j).
      \notag
    \end{align}
    This is the gauge normalization used in
    \cite[Appendix~C.4, lines~2025--2029]{Cirac2016MPDO_arXiv}.
\end{theorem}

\begin{proof}\leanok
    Apply the Figure~8 comparison to the active corner and the distinguished
    blocked reference corner. Both coefficients are positive, so horizontal
    canonicality and positivity identify their Gram-dressed normal tensors.
    Normal-tensor rigidity then makes $X_j^\dagger X_j$ a positive real
    multiple of the identity. Division by the positive square root gives a
    unitary matrix.
\end{proof}

\begin{theorem}[Literal unitary normalization of active product gauges]
    \label{thm:mpdo_literal_cpsv_active_product_gauge_unitary_normalization}
    \lean{MPOTensor.RetainedProductSpectralFamily.FlatBlockedBNTComparison.%
exists_unitaryNormalization_of_cpsvCanonicalForm}
    \leanok
    \uses{def:mpdo, def:cpsv_canonical_form,
      thm:mpdo_active_product_corner_blocked_bnt_coverage,
      thm:mpdo_blocked_reference_corners,
      thm:mpdo_literal_cpsv_active_product_coefficient_phase_pos,
      thm:cpsv_literal_positive_relative_gram_scalar}
    Under the hypotheses of the preceding literal positivity theorem, for
    every active product corner $j$ there is a real number $\omega_j>0$ such
    that
    \begin{align}
      X_j^\dagger X_j
      &=\omega_j\Id,
      \notag\\
      \omega_j^{-1/2}X_j
      &\in\mathrm{U}(D_j).
      \notag
    \end{align}
\end{theorem}

\begin{proof}\leanok
    Compress the blocked vertical tensor to the active corner and to the
    distinguished blocked reference corner. Their coefficients are positive
    by the preceding theorem and by positivity of the blocked multiplicity
    weight. The literal Figure~8 comparison gives
    $X_j^\dagger X_j=\omega_j\Id$ with $\omega_j>0$; division by
    $\sqrt{\omega_j}$ gives a unitary matrix.
\end{proof}

\begin{definition}[Original-label corner family]
    \label{def:mpdo_original_corner_family}
    \lean{MPOTensor.RetainedProductSpectralFamily.OriginalCornerFamily}
    \leanok
    \uses{def:mpdo_retained_product_spectral_family}
    An original-label corner family for the retained product records, for
    every active corner of each copy pair, an original one-site BNT label, a
    positive coefficient, and an isometric inclusion into the raw product
    bond space. Distinct corners belonging to one copy pair have orthogonal
    ranges. Each inclusion intertwines the raw product tensor with the stated
    positive multiple of its original BNT representative, and the resulting
    corners reconstruct every tensor letter exactly.

    \textbf{Local fix (Figure-11 fixed-pair support):} A copy pair with no
    active corner is represented by an empty family; no unsupported corner is
    inserted. See
    \path{docs/paper-gaps/cpsv16_figure11_per_pair_support.tex}.
\end{definition}

\begin{theorem}[Positive original-label corner decomposition]
    \label{thm:mpdo_original_corner_family_exists}
    \lean{MPOTensor.RetainedProductSpectralFamily.exists_originalCornerFamily}
    \leanok
    \uses{def:mpdo_original_corner_family,
      thm:mpdo_active_product_corner_blocked_bnt_coverage,
      thm:mpdo_transported_vertical_sector_coefficient_comparison,
      lem:mpdo_normalized_vertical_sector_positive_trace,
      thm:mpdo_active_product_coefficient_phase_pos,
      thm:mpdo_active_product_gauge_unitary_normalization,
      thm:mpdo_normalized_grouped_sector_maps}
    Let the active product corners be covered by blocked BNT representatives,
    let the blocked representatives be identified with the original BNT by
    the trace-ratio unitary conjugacies, and suppose that the active
    gauge--phase coefficients are positive. If
    $X_j^\dagger X_j=\omega_j\Id$ with $\omega_j>0$, then every retained copy
    pair has an original-label corner family. Its coefficient at an active
    corner is
    \begin{align}
      \frac{\lambda_j\zeta_j}{w_{\alpha,q}w_{\beta,r}}
      \frac{\tr D_{\gamma(j)}}
           {\tr D^{(2)}_{\sigma(\gamma(j))}},
      \notag
    \end{align}
    which is positive. The normalized inclusions are isometries with
    pairwise orthogonal ranges, intertwine the raw product tensor with these
    positive multiples of the original BNT representatives, and reconstruct
    every raw product letter exactly.

    \textbf{Local fix (Figure-11 fixed-pair support):} Empty active families
    remain empty. Exact weighted reconstruction and positivity of the outer
    copy weight imply that the corresponding raw product vanishes. See
    \path{docs/paper-gaps/cpsv16_figure11_per_pair_support.tex}.
\end{theorem}

\begin{proof}\leanok
    Transport each local inclusion first by the normalized active comparison
    gauge and then by the unitary identifying the blocked representative with
    its original one-site label. The first normalization preserves isometry
    and orthogonality; the second unitary transport does likewise. The two
    conjugacy equations give the stated coefficient and intertwining identity.
    Positivity follows from the active gauge--phase coefficient, the two
    positive multiplicity traces, and the positive outer copy weight.
    Conjugating each corner term through the two transports recovers its local
    spectral term. Summing these identities gives the weighted
    reconstruction, and cancellation of the nonzero outer weight gives the
    raw reconstruction.
\end{proof}

\begin{theorem}[Fusion coisometries from original-label corners]
    \label{thm:mpdo_original_corner_family_to_fusion_coisometry}
    \lean{MPOTensor.RetainedProductSpectralFamily.OriginalCornerFamily.%
toBNTFusionCoisometryFamily}
    \leanok
    \uses{def:mpdo_original_corner_family,
      def:mpdo_bnt_fusion_coisometry_family,
      thm:vertical_sector_coisometry}
    Given an original-label corner family, choose the distinguished retained
    copy of every BNT label. For fixed
    labels $\alpha,\beta$, let the multiplicity of $\gamma$ be the number of
    active corners of the distinguished copy pair whose original label is
    $\gamma$, and put the corresponding positive corner coefficients on the
    diagonal of $\chi_{\alpha,\beta,\gamma}$. The block row of the adjoints
    of the corner inclusions, regrouped by these label fibers, defines a
    fusion coisometry. It satisfies
    \begin{align}
      \chi_{\alpha,\beta,\gamma;r}
      &>0,
      \notag\\
      U_{\alpha,\beta}U_{\alpha,\beta}^\dagger
      &=1,
      \notag\\
      U_{\alpha,\beta}(M_\alpha M_\beta)^{ij}
        U_{\alpha,\beta}^\dagger
      &=\bigoplus_\gamma
          \chi_{\alpha,\beta,\gamma}\otimes M_\gamma^{ij},
      \notag
    \end{align}
    together with the exact reconstruction
    \begin{align}
      (M_\alpha M_\beta)^{ij}
      &=U_{\alpha,\beta}^\dagger
          \left(\bigoplus_\gamma
            \chi_{\alpha,\beta,\gamma}\otimes M_\gamma^{ij}\right)
          U_{\alpha,\beta}.
      \notag
    \end{align}
    Empty label fibers contribute zero-dimensional diagonal blocks.
    This is the decomposition of CPSV16, Appendix C.4, lines 2020--2029.

    \textbf{Scope restriction (active product BNT):} Only active product
    corners are retained. A label absent from the distinguished copy pair
    has zero fusion multiplicity. See
    \path{docs/paper-gaps/cpsv16_bnt_uniqueness_zero_coefficient.tex}.

    \textbf{Local fix (Figure-11 fixed-pair support):} The distinguished copy
    pair may have an empty active family; no unsupported corner is inserted.
    See \path{docs/paper-gaps/cpsv16_figure11_per_pair_support.tex}.

    \textbf{Local fix (Figure-11 fusion coisometry):} The retained-row map is
    a coisometry onto the active direct sum. Exact reconstruction includes
    any common zero corner discarded by the forward map. See
    \path{docs/paper-gaps/cpsv16_figure11_fusion_coisometry.tex}.
\end{theorem}

\begin{proof}\leanok
    For each distinguished copy pair, form the block row from the adjoints of
    its pairwise orthogonal isometric corner inclusions. The block-row
    identity makes this map a coisometry. Its forward conjugation is block
    diagonal over the active corners, while the corner reconstruction gives
    the reverse identity. Regrouping the active-corner sum by the original
    label turns each label fiber into the diagonal matrix whose entries are
    precisely the positive corner coefficients. These identities are valid
    without separate cases when a fiber or the whole active family is empty.
\end{proof}

\begin{theorem}[Fusion coisometries from normalized product corners]
    \label{thm:mpdo_normalized_product_corners_to_fusion_coisometry}
    \lean{MPOTensor.RetainedProductSpectralFamily.%
exists_bntFusionCoisometryFamily}
    \leanok
    \uses{thm:mpdo_original_corner_family_exists,
      thm:mpdo_original_corner_family_to_fusion_coisometry}
    Suppose that every active product corner has positive gauge--phase
    coefficient and that its gauge Gram matrix is a positive scalar multiple
    of the identity. Given the trace-ratio unitary correspondence between the
    one-site and two-site vertical sectors, there are positive diagonal
    matrices $\chi_{\alpha,\beta,\gamma}$ and coisometries
    $U_{\alpha,\beta}$ satisfying the forward fusion identity and exact
    reconstruction for every product tensor $M_\alpha M_\beta$. This is the
    active-corner construction in
    \cite[Appendix~C.4, lines~2020--2029]{Cirac2016MPDO_arXiv}.
\end{theorem}

\begin{proof}\leanok
    The positive coefficients and normalized gauges give an original-label
    corner family. Apply the preceding block-row construction to its
    orthogonal isometric inclusions.
\end{proof}

\begin{theorem}[Positive fusion from transported vertical sectors]
    \label{thm:mpdo_transported_vertical_product_positive_fusion_decomposition}
    \lean{MPOTensor.transportedVerticalSector_exists_positiveFusionDecomposition}
    \leanok
    \uses{lem:cpsv_bnt_block_dimensions_positive,
      thm:mpdo_transported_vertical_sector_coefficient_comparison,
      thm:mpdo_retained_product_spectral_family_exists,
      thm:mpdo_active_product_corner_blocked_bnt_coverage,
      thm:mpdo_active_product_coefficient_phase_pos,
      thm:mpdo_active_product_gauge_unitary_normalization,
      thm:mpdo_original_corner_family_exists,
      thm:mpdo_original_corner_family_to_fusion_coisometry,
      thm:mpdo_normalized_product_corners_to_fusion_coisometry}
    Under the one-site and two-site vertical basis-of-normal-tensors
    decompositions and the two completely positive, trace-preserving
    physical-closure transport identities, assume that the one-site tensor
    is in normalized BNT-refined horizontal form and defines a matrix product
    density operator. Then there are positive diagonal multiplicity matrices
    $\chi_{\alpha,\beta,\gamma}$ and matrices $U_{\alpha,\beta}$ such that
    every diagonal entry is positive and
    \begin{align}
      U_{\alpha,\beta}U_{\alpha,\beta}^{\dagger}
      &=1,
      \notag\\
      U_{\alpha,\beta}(M_\alpha M_\beta)^{ij}
        U_{\alpha,\beta}^{\dagger}
      &=\bigoplus_\gamma
          \chi_{\alpha,\beta,\gamma}\otimes M_\gamma^{ij},
      \notag
    \end{align}
    while the exact reconstruction is
    \begin{align}
      (M_\alpha M_\beta)^{ij}
      &=U_{\alpha,\beta}^{\dagger}
          \left(\bigoplus_\gamma
            \chi_{\alpha,\beta,\gamma}\otimes M_\gamma^{ij}\right)
          U_{\alpha,\beta}.
      \notag
    \end{align}
    This is the positive fusion decomposition of CPSV16, Appendix C.4,
    lines 2020--2029.

    \textbf{Scope restriction (active product BNT):} Only active product
    corners occur. A label absent from a fixed product pair has zero fusion
    multiplicity. See
    \path{docs/paper-gaps/cpsv16_bnt_uniqueness_zero_coefficient.tex}.

    \textbf{Local fix (Figure-11 fixed-pair support):} A fixed product pair
    may have no active corner, and no unsupported sector is inserted. See
    \path{docs/paper-gaps/cpsv16_figure11_per_pair_support.tex}.

    \textbf{Local fix (Figure-11 fusion coisometry):} The retained-row map is
    a coisometry onto the active direct sum, and its adjoint gives the exact
    reconstruction. See
    \path{docs/paper-gaps/cpsv16_figure11_fusion_coisometry.tex}.
\end{theorem}

\begin{proof}\leanok
    First decompose every retained copy-pair tensor into its active normal
    corners and compare these corners with the two-site vertical basis.
    Positivity of matrix product density operators makes every resulting
    scalar positive, while the Gram comparison normalizes each gauge to a
    unitary. The inverse vertical-sector transport then returns the corners
    to the original one-site labels. Their orthogonal isometric inclusions
    form the rows of $U_{\alpha,\beta}$. Grouping these rows by the original
    label gives the positive diagonal matrices $\chi_{\alpha,\beta,\gamma}$;
    corner intertwining and corner reconstruction give the two displayed
    identities.
\end{proof}

\begin{theorem}[Positive fusion from literal CPSV canonical form]
    \label{thm:mpdo_literal_cpsv_vertical_product_positive_fusion_decomposition}
    \lean{MPOTensor.%
exists_positiveFusionDecomposition_of_unitaryBlockEquiv_of_cpsvCanonicalForm}
    \leanok
    \uses{def:mpdo, def:cpsv_canonical_form,
      thm:mpdo_transported_vertical_sector_coefficient_comparison,
      thm:mpdo_literal_cpsv_retained_product_spectral_family,
      thm:mpdo_active_product_corner_blocked_bnt_coverage,
      thm:mpdo_literal_cpsv_active_product_coefficient_phase_pos,
      thm:mpdo_literal_cpsv_active_product_gauge_unitary_normalization,
      thm:mpdo_original_corner_family_exists,
      thm:mpdo_original_corner_family_to_fusion_coisometry,
      thm:mpdo_normalized_product_corners_to_fusion_coisometry}
    Let $M$ generate matrix product density operators, and suppose that its
    doubled-index MPS tensor is in literal CPSV canonical form. Given one-site
    and two-site vertical canonical decompositions and their trace-ratio
    unitary sector correspondence, there are positive diagonal matrices
    $\chi_{\alpha,\beta,\gamma}$ and matrices $U_{\alpha,\beta}$ such that
    \begin{align}
      U_{\alpha,\beta}U_{\alpha,\beta}^{\dagger}
      &=1,
      \notag\\
      U_{\alpha,\beta}(M_\alpha M_\beta)^{ij}
        U_{\alpha,\beta}^{\dagger}
      &=\bigoplus_\gamma
        \chi_{\alpha,\beta,\gamma}\otimes M_\gamma^{ij},
      \notag\\
      (M_\alpha M_\beta)^{ij}
      &=U_{\alpha,\beta}^{\dagger}
        \left(\bigoplus_\gamma
          \chi_{\alpha,\beta,\gamma}\otimes M_\gamma^{ij}\right)
        U_{\alpha,\beta}.
      \notag
    \end{align}
    Every diagonal entry of $\chi_{\alpha,\beta,\gamma}$ is positive.
    This is the positive fusion decomposition in
    \cite[Appendix~C.4, lines~2020--2029]{Cirac2016MPDO_arXiv}.

    \textbf{Scope restriction (active product BNT):} Only active product
    corners occur. A label absent from a fixed product pair has zero fusion
    multiplicity. See
    \path{docs/paper-gaps/cpsv16_bnt_uniqueness_zero_coefficient.tex}.

    \textbf{Local fix (Figure-11 fixed-pair support):} A fixed product pair
    may have no active corner, and no unsupported sector is inserted. See
    \path{docs/paper-gaps/cpsv16_figure11_per_pair_support.tex}.

    \textbf{Local fix (Figure-11 fusion coisometry):} The retained-row map is
    a coisometry onto the active direct sum, and its adjoint gives the exact
    reconstruction. See
    \path{docs/paper-gaps/cpsv16_figure11_fusion_coisometry.tex}.
\end{theorem}

\begin{proof}\leanok
    Literal canonical form gives simultaneous spectra for all retained product
    corners. Compare each active corner with the blocked vertical BNT. Literal
    sector separation makes its coefficient positive, and the literal
    Figure~8 comparison normalizes its gauge to a unitary. Transporting these
    corners through the one-site--two-site sector correspondence gives
    positive original-label corners; their adjoints form the rows of the
    coisometries $U_{\alpha,\beta}$ and yield both displayed identities.
\end{proof}

\begin{theorem}[Idempotent trace scalars from blocked representations]
    \label{thm:mpdo_blocked_representations_idempotent_trace_scalars}
    \lean{MPOTensor.hasIdempotentCoefficientForm_of_blockedRepresentations}
    \leanok
    \uses{def:mpdo_bnt_label_idempotent_coefficient_form,
      def:mpdo_bnt_fusion_coisometry_family,
      def:mpdo_bnt_eventual_power_sums_multiplicity_comparison,
      thm:mpdo_blocked_vertical_operator_representations}
    Fix one-site and two-site vertical canonical decompositions whose sectors
    are identified by the transported unitary relabelling. Suppose that the
    product tensors have positive diagonal fusion matrices
    $\chi_{\alpha,\beta,\gamma}$ and exact forward and reverse coisometric
    fusion identities. If the blocked operator has the two simultaneous
    representations of
    Theorem~\ref{thm:mpdo_blocked_vertical_operator_representations}, then
    \begin{align}
      m_\gamma
      &=\sum_{\alpha,\beta}
        \tr(\chi_{\alpha,\beta,\gamma})m_\alpha m_\beta.
      \notag
    \end{align}
    This is the idempotent coefficient identity in
    \cite[Appendix~C.4, lines~2030--2042]{Cirac2016MPDO_arXiv}.
\end{theorem}

\begin{proof}\leanok
    For every sufficiently large $L$, expand the product of one-site sector
    operators with the fusion identity. Eventual linear independence of the
    BNT operators and the two blocked representations then identify, for
    every $\gamma$, the two positive power sums
    \begin{align}
      \tr\!\left[
        \left(\frac{m_\gamma}{n_{\sigma(\gamma)}}
          \nu_{\sigma(\gamma)}\right)^L\right]
      &=
      \sum_{\alpha,\beta}
        \tr\!\left[
          \left(\mu_\alpha\otimes\mu_\beta\otimes
            \chi_{\alpha,\beta,\gamma}\right)^L\right].
      \notag
    \end{align}
    The eventual power-sum comparison identifies the two multisets of positive
    entries. Summing this multiset equality and cancelling
    $n_{\sigma(\gamma)}$ gives the displayed length-one identity.
\end{proof}

\begin{theorem}[Positive fusion decomposition of vertical products]
    \label{thm:mpdo_vertical_product_positive_fusion_decomposition}
    \lean{MPOTensor.exists_positiveFusionDecomposition_of_isRFPViaTS}
    \leanok
    \uses{def:is_rfp_via_ts,
      thm:vertical_cf_of_horizontal_cf,
      thm:mpdo_transported_vertical_product_positive_fusion_decomposition}
    Let $M$ be in normalized BNT-refined horizontal form, generate matrix
    product density operators, and satisfy the renormalization fixed-point condition of
    Definition~\ref{def:is_rfp_via_ts}. Then one may choose a vertical basis
    of normal tensors $\{M_\alpha\}_\alpha$ for which there are positive
    diagonal matrices $\chi_{\alpha,\beta,\gamma}$ and matrices
    $U_{\alpha,\beta}$. Every diagonal entry of
    $\chi_{\alpha,\beta,\gamma}$ is positive, and
    \begin{align}
      U_{\alpha,\beta}U_{\alpha,\beta}^{\dagger}
      &=1,
      \notag\\
      U_{\alpha,\beta}(M_\alpha M_\beta)^{ij}
        U_{\alpha,\beta}^{\dagger}
      &=\bigoplus_\gamma
          \chi_{\alpha,\beta,\gamma}\otimes M_\gamma^{ij},
      \notag\\
      (M_\alpha M_\beta)^{ij}
      &=U_{\alpha,\beta}^{\dagger}
          \left(\bigoplus_\gamma
            \chi_{\alpha,\beta,\gamma}\otimes M_\gamma^{ij}\right)
          U_{\alpha,\beta}.
      \notag
    \end{align}
    Thus the renormalization fixed-point and matrix-product-density-operator
    assumptions, together with normalized BNT-refined horizontal form, imply
    the positive fusion decomposition of CPSV16, Appendix C.4, lines
    2020--2029.

    \textbf{Scope restriction (active product BNT):} Only active product
    corners occur. A label absent from a fixed product pair has zero fusion
    multiplicity. See
    \path{docs/paper-gaps/cpsv16_bnt_uniqueness_zero_coefficient.tex}.

    \textbf{Local fix (Figure-11 fixed-pair support):} A fixed product pair
    may have no active corner, and no unsupported sector is inserted. See
    \path{docs/paper-gaps/cpsv16_figure11_per_pair_support.tex}.

    \textbf{Local fix (Figure-11 fusion coisometry):} The retained-row map is
    a coisometry onto the active direct sum, and its adjoint gives the exact
    reconstruction. See
    \path{docs/paper-gaps/cpsv16_figure11_fusion_coisometry.tex}.
\end{theorem}

\begin{proof}\leanok
    Apply the proved normalized BNT-refined specialization of
    Proposition~4.13 to $M$ and to its two-site blocking. In each case, retain
    the CPSV basis-of-normal-tensors assertion, the positive multiplicity
    matrices, and both coisometric decomposition identities.
    Unpack the two completely positive, trace-preserving maps in the
    renormalization fixed-point condition. The preceding transported-sector
    theorem then supplies the positive diagonal family, the row
    coisometries, the forward fusion identity, and exact reconstruction for
    the chosen one-site vertical basis.
\end{proof}

\begin{theorem}[Two-site closure of the two-site blocking]
    \label{thm:mpdo_two_site_closure_two_site_blocking}
    \lean{MPOTensor.physClose2_blockTwo_eq_physClose4}
    \leanok
    \uses{def:phys_close, def:mpdo_two_site_blocking,
      def:mpdo_four_site_physical_closure}
    After identifying two blocked physical indices with four original
    indices, one has $(K^{[2]})_2(X)=K_4(X)$.
\end{theorem}

\begin{proof}\leanok
    Expand the two blocked tensor matrices and reassociate their product.
\end{proof}
